% CMPE 220 — Assignment 1: Software CPU Design (report)
% Author: Faiq Malik  |  Spring 2026
%
% Figures for \includegraphics{...} live in ../figures/ (compile from reports/).
%
% Build PDF (from repository root, optional output directory):
%   pdflatex -interaction=nonstopmode -output-directory=reports reports/Assignment1_Software_CPU_Report.tex
%   pdflatex -interaction=nonstopmode -output-directory=reports reports/Assignment1_Software_CPU_Report.tex
%
% Edit \githuburl below before submitting.

\documentclass[12pt]{article}

\usepackage[margin=1in]{geometry}
\usepackage[doublespacing]{setspace}
\usepackage{url}

\pagestyle{plain}

\newcommand{\githuburl}{https://github.com/USERNAME/REPONAME}

\begin{document}

% ---------------------------------------------------------------------------
% Title page
% ---------------------------------------------------------------------------
\begin{titlepage}
    \thispagestyle{empty}
    \centering
    \vspace*{1.5in}

    {\LARGE\bfseries CMPE 220 --- System Software}\\[0.35in]
    {\Large Spring 2026}\\[0.75in]

    {\LARGE\bfseries Assignment 1}\\[0.2in]
    {\LARGE\bfseries Software CPU Design and Implementation}\\[1.25in]

    {\large\bfseries Team members}\\[0.2in]
    {\large Faiq Malik}\\[1.5in]

    {\large Date: \today}
\end{titlepage}

\clearpage
\setcounter{page}{1}

% ---------------------------------------------------------------------------
% GitHub repository
% ---------------------------------------------------------------------------
\section*{GitHub Repository}

\noindent
The complete source code, any built artifacts or example binaries included in the repository, demonstration video, and documentation are maintained at:

\vspace{0.5in}
\begin{center}
\large\url{\githuburl}
\end{center}

\vspace{0.35in}
\noindent
The repository root contains \texttt{cpu16\_core.cpp} (single-file CPU, assembler, and emulator), demonstration media under \texttt{demos/} (including \texttt{demos/Demo.mp4}), diagrams under \texttt{figures/} when present, any prebuilt program images (for example \texttt{fib.bin}), and this report under \texttt{reports/}.

\clearpage

% ---------------------------------------------------------------------------
% Download, compile, and run
% ---------------------------------------------------------------------------
\section*{Download, Compile, and Run}

\subsection*{Download}

\noindent
Clone the repository with Git, or download a ZIP archive from the GitHub page above and extract it to a working directory on macOS or Linux.

\subsection*{Requirements}

\begin{itemize}
  \item C++ compiler with C++17 support (\texttt{g++} recommended).
  \item Terminal access.
\end{itemize}

\subsection*{Compile}

\noindent
From the repository root directory:

\begin{verbatim}
g++ -std=c++17 -O2 -o cpu16 cpu16_core.cpp
\end{verbatim}

\noindent
This produces the executable \texttt{cpu16}, which implements assembler mode (\texttt{asm}), bare-memory emulator mode (\texttt{emu}), and assemble-and-run mode (\texttt{run}).

\subsection*{Run (examples)}

\noindent
Typical commands when example assembly sources are under \texttt{examples/}:

\begin{verbatim}
./cpu16 run examples/hello.asm
./cpu16 run examples/timer.asm
./cpu16 asm examples/fib.asm -o fib.bin
./cpu16 emu fib.bin --base 0x0000 --pc 0x0100 --dump 0x0100 0x0140
\end{verbatim}

\noindent
The Fibonacci listing uses \texttt{.org 0x0100}; when emulating \texttt{fib.bin}, set the program counter to \texttt{0x0100} as in the command above. If \texttt{fib.bin} lives in the repository root instead of beside the assembly file, use the same arguments with that path, for example:

\begin{verbatim}
./cpu16 emu fib.bin --base 0x0000 --pc 0x0100 --dump 0x0100 0x0140
\end{verbatim}

\noindent
To assemble factorial from source when \texttt{examples/fact.asm} is available, then load and dump memory:

\begin{verbatim}
./cpu16 asm examples/fact.asm -o fact.bin
./cpu16 emu fact.bin --base 0x0000 --pc 0x0000 --dump 0x0000 0x01FF
\end{verbatim}

\subsection*{Video}

\noindent
The Fibonacci demonstration required for the assignment is included as \texttt{demos/Demo.mp4} in the repository.

\clearpage

% ---------------------------------------------------------------------------
% Contributions
% ---------------------------------------------------------------------------
\section*{Team Member Contributions}

\subsection*{Faiq Malik}

\noindent
Contributions include:

\begin{itemize}
  \item \textbf{Architecture and ISA:} 16-bit CPU organization (registers, ALU, control flow, memory map, MMIO), instruction formats and encodings, and flag semantics as implemented in \texttt{cpu16\_core.cpp}.
  \item \textbf{Emulator:} Fetch--decode--execute, register file, ALU with condition flags, program counter, stack conventions, per-instruction timer step, and UART/timer MMIO behavior.
  \item \textbf{Assembler:} Two-pass assembly with labels, literals, directives (\texttt{.org}, \texttt{.word}, \texttt{.stringz}), and resolution of forward references.
  \item \textbf{Example programs:} Hello world, timer, Fibonacci, and factorial (or equivalent) assembly used for demonstration and grading.
  \item \textbf{Deliverables:} Repository layout, grader-facing build and run documentation (including this report), demonstration video (\texttt{demos/Demo.mp4}), supporting assets as required, and this report submitted as PDF for the course.
\end{itemize}

\end{document}
